% Actual class: 02
% Slide decks: 2_regression_classification; L2_summary
% Exact scope: 2_regression_classification pp. 1--65; L2_summary pp. 1--23
% Video supplement: Lecture2_Part1 and Lecture2_Part2, complete visual scope checked
% Human verified: Yes

\section{第 02 节课：Regression 与 Classification}
\label{lecture:SC4001-L02}

\lecturemeta
  {SC4001-L02}
  {2026-08-21}
  {2\_regression\_classification；L2\_summary}
  {主讲义 pp.~1--65；Summary pp.~1--23}
  {Lecture2\_Part1（40:39）与 Lecture2\_Part2（35:37）；关键画面与讲义进度已核对}
  {已确认（2026-08-21；检查题完成并订正）}

\subsection{本讲主线}

Predictive AI（预测式人工智能）从 features（特征）$\bm x$ 预测 labels（标签）$y$。Continuous labels（连续标签）对应 regression（回归），discrete labels（离散标签）对应 classification（分类）；二者共享 affine score（仿射得分）$u=\bm w^{\mathsf T}\bm x+b$，但使用不同的 output interpretation（输出含义）与 loss function（损失函数）。

\begin{figure}[htbp]
  \centering
  \resizebox{0.94\textwidth}{!}{% Original diagram based on the relationships in Lecture 2; no slide artwork copied.
\begin{tikzpicture}[
  >=Latex,
  block/.style={draw=noteBlue,very thick,rounded corners=3pt,fill=blue!7,minimum width=2.65cm,minimum height=0.85cm,align=center},
  output/.style={draw=ntuRed,very thick,rounded corners=3pt,fill=red!6,minimum width=3.5cm,minimum height=0.95cm,align=center},
  flow/.style={->,thick,draw=noteBlue},
  label/.style={font=\scriptsize,align=center,fill=white,inner sep=1.5pt}
]
  \node[block] (x) at (0,0) {Features $\bm x$};
  \node[block] (u) at (4.0,0) {$u=\bm w^{\mathsf T}\bm x+b$};
  \node[output] (reg) at (8.5,1.25) {Regression\\$y=f(u)\in\mathbb R$};
  \node[output] (cls) at (8.5,-1.25) {Classification\\$p=\sigma(u)$，$\hat y=\mathbf 1(p>0.5)$};
  \draw[flow] (x) -- node[label,above]{affine computation} (u);
  \draw[flow] (u) -- node[label,above,sloped]{continuous target} (reg);
  \draw[flow] (u) -- node[label,below,sloped]{binary target} (cls);
\end{tikzpicture}
}
  \caption{Regression 与 binary classification（二分类）共享 affine computation（仿射计算），但输出与损失的含义不同。}
  \label{fig:sc4001-l02-task-map}
\end{figure}

\lecturetopic{SC4001-L02-T01}{Regression、Classification 与课程术语}

\begin{center}
\small
\begin{tabularx}{\textwidth}{@{}>{\raggedright\arraybackslash}p{2.5cm}>{\raggedright\arraybackslash}p{2.7cm}>{\raggedright\arraybackslash}p{3.2cm}X@{}}
  \toprule
  Task & Target & 典型 output unit & 本讲 objective（目标函数） \\
  \midrule
  Regression（回归） & $d\in\mathbb R$ & Linear neuron；continuous sigmoid unit & Square error（平方误差） \\
  Binary classification（二分类） & $d\in\{0,1\}$ & Discrete perceptron；logistic regression neuron & Binary cross-entropy（二元交叉熵） \\
  \bottomrule
\end{tabularx}
\end{center}

课程把使用 sigmoid activation（Sigmoid 激活）的连续 unit 也称为 perceptron；一般文献常把 perceptron 专指 threshold unit（阈值单元）。本笔记沿用课程术语，并用 discrete perceptron（离散感知机）明确指 threshold classifier（阈值分类器）。

\lecturetopic{SC4001-L02-T02}{Square-error Gradient、SGD 与 GD}

对一个 training sample（训练样本）$(\bm x,d)$，令
\[
  u=\bm w^{\mathsf T}\bm x+b,\qquad y=f(u),\qquad
  J=\frac12(d-y)^2.
\]
Chain rule（链式法则）给出
\[
  \frac{\partial J}{\partial u}=-(d-y)f'(u),\qquad
  \boxed{\nabla_{\bm w}J=-(d-y)f'(u)\bm x},\qquad
  \boxed{\frac{\partial J}{\partial b}=-(d-y)f'(u)}.
\]
因此课程定义的 SGD（每个 sample 更新一次）为
\[
  \boxed{\bm w\leftarrow\bm w+\alpha(d-y)f'(u)\bm x},\qquad
  \boxed{b\leftarrow b+\alpha(d-y)f'(u)}.
\]

若 $P$ 个 samples 按行组成 $\bm X\in\mathbb R^{P\times n}$，且
\[
  \bm u=\bm X\bm w+b\bm 1_P,\qquad \bm y=f(\bm u),
\]
则 full-batch GD（全批量梯度下降）的更新是
\[
  \boxed{\bm w\leftarrow\bm w+
    \alpha\bm X^{\mathsf T}\!\left[(\bm d-\bm y)\odot f'(\bm u)\right]},
  \qquad
  \boxed{b\leftarrow b+
    \alpha\bm 1_P^{\mathsf T}\!\left[(\bm d-\bm y)\odot f'(\bm u)\right]}.
\]
$\odot$ 表示 element-wise product（逐元素乘法）。讲义使用 sum loss（求和损失）；若实现改用 mean loss（平均损失），gradient scale（梯度尺度）和适用的 learning rate（学习率）也会改变。

\lecturetopic{SC4001-L02-T03}{Linear 与 Sigmoid Regression}

Linear neuron（线性神经元）使用 $f(u)=u$、$f'(u)=1$：
\[
  y=\bm w^{\mathsf T}\bm x+b,\qquad
  \bm w\leftarrow\bm w+\alpha(d-y)\bm x,\qquad
  b\leftarrow b+\alpha(d-y).
\]
它学习 feature space（特征空间）中的 best-fit hyperplane（最佳拟合超平面）；只有数据可被精确线性拟合时，该超平面才通过所有 training points（训练点）。

Continuous sigmoid unit（连续 Sigmoid 单元）使用
\[
  \sigma(u)=\frac{1}{1+e^{-u}},\qquad
  \sigma'(u)=\sigma(u)[1-\sigma(u)].
\]
若 target range（目标值域）不是 $(0,1)$，应使用 scaled activation（缩放激活）并重新求导。例如 $f(u)=4\sigma(u)-1$ 的值域为 $(-1,3)$，且
\[
  f'(u)=4\sigma(u)[1-\sigma(u)],
\]
不能把 transformed output（变换后的输出）直接代入 $y(1-y)$。单个 sigmoid unit 虽为 nonlinear mapping（非线性映射），但只能表示 $f(\bm w^{\mathsf T}\bm x+b)$ 这一受限函数族。

常用 derivative（导数）还包括
\[
  \frac{d}{du}\tanh u=1-\tanh^2u.
\]

\textbf{对应例题：}\relatedexercise{SC4001-L02-E01}{Lecture Example 1：Scaled Sigmoid Regression}。

\lecturetopic{SC4001-L02-T04}{Decision Boundary 与 Logistic Regression}

Linear discriminant function（线性判别函数）为
\[
  g(\bm x)=\bm w^{\mathsf T}\bm x+b.
\]
Decision boundary（决策边界）满足 $g(\bm x)=0$；$\bm w$ 是 boundary 的 normal vector（法向量），$b$ 控制位置。可约定 $g(\bm x)>0$ 为 class 1，$g(\bm x)\le0$ 为 class 0。

\begin{figure}[htbp]
  \centering
  \resizebox{0.78\textwidth}{!}{% Original geometric illustration; coordinates are illustrative.
\begin{tikzpicture}[>=Latex,scale=1.0]
  \draw[->,thick] (-3.2,0) -- (3.3,0) node[right] {$x_1$};
  \draw[->,thick] (0,-2.2) -- (0,3.0) node[above] {$x_2$};
  \draw[very thick,draw=ntuRed] (-3,2.6) -- (3,-2.2)
    node[above right,font=\small] {$g(\bm x)=0$};

  \foreach \p in {(-2.4,0.5),(-1.9,0.0),(-1.4,0.8),(-0.9,-0.3),(-0.4,0.1)}
    \fill[noteBlue] \p circle (2.8pt);
  \foreach \p in {(0.5,1.6),(1.0,2.1),(1.5,1.3),(2.0,2.4),(2.5,1.7)}
    \draw[ntuRed,very thick] \p ++(-0.11,-0.11) -- ++(0.22,0.22)
      \p ++(-0.11,0.11) -- ++(0.22,-0.22);

  \node[noteBlue,font=\small] at (-2.0,-1.15) {$g(\bm x)<0$};
  \node[ntuRed,font=\small] at (1.9,0.65) {$g(\bm x)>0$};
  \coordinate (q) at (0,0.2);
  \draw[->,very thick,draw=verifyOrange] (q) -- (0.8,1.2)
    node[right,font=\small] {$\bm w$};
  \fill[verifyOrange] (q) circle (2pt);
\end{tikzpicture}
}
  \caption{Linear decision boundary（线性决策边界）：边界两侧由 $g(\bm x)$ 的符号区分，$\bm w$ 与边界垂直。}
  \label{fig:sc4001-l02-decision-boundary}
\end{figure}

Logistic regression neuron（逻辑回归神经元）令
\[
  p=f(u)=\sigma(u)=P(y=1\mid\bm x),\qquad
  \hat y=\mathbf 1(p>0.5).
\]
必须区分 probability（概率）$p$ 与 predicted label（预测标签）$\hat y$。因为 sigmoid strictly increasing（严格递增）且 $\sigma(0)=0.5$，
\[
  p>0.5\iff u>0,
\]
所以其 decision boundary 仍是 $\bm w^{\mathsf T}\bm x+b=0$。课程采用 strict threshold（严格阈值），故 $u=0$ 时 $p=0.5$、$\hat y=0$。

\FloatBarrier
\lecturetopic{SC4001-L02-T05}{Binary Cross-Entropy 与 Logistic Gradient}

对 $d\in\{0,1\}$，binary cross-entropy（BCE）为
\[
  J=-d\ln p-(1-d)\ln(1-p),\qquad p=\sigma(u).
\]
利用
\[
  \frac{\partial J}{\partial p}=\frac{p-d}{p(1-p)},\qquad
  \frac{\partial p}{\partial u}=p(1-p),
\]
可得简洁的 logit gradient（对激活前值的梯度）：
\[
  \boxed{\frac{\partial J}{\partial u}=p-d}.
\]
于是 single-sample update（单样本更新）为
\[
  \boxed{\bm w\leftarrow\bm w+\alpha(d-p)\bm x},\qquad
  \boxed{b\leftarrow b+\alpha(d-p)},
\]
batch update（批量更新）为
\[
  \boxed{\bm w\leftarrow\bm w+\alpha\bm X^{\mathsf T}[\bm d-\sigma(\bm u)]},
  \qquad
  \boxed{b\leftarrow b+\alpha\bm 1_P^{\mathsf T}[\bm d-\sigma(\bm u)]}.
\]
BCE 是 differentiable training objective（可微训练目标）；classification error（分类错误数）$\sum_p\mathbf 1(d_p\ne\hat y_p)$ 适合 evaluation（评估），但不适合直接做 gradient-based optimization（基于梯度的优化）。

\textbf{对应例题：}\relatedexercise{SC4001-L02-E02}{Lecture Example 2：GD for Logistic Regression}。

\lecturetopic{SC4001-L02-T06}{Linear Separability、Learning Rate 与 Autograd}

Single logistic neuron 的 nonlinearity（非线性）改变 probability 随 $u$ 的变化，却不改变 threshold condition（阈值条件）$u=0$；因此它只能处理 linearly separable patterns（线性可分模式），不能直接解决 XOR。

\begin{figure}[htbp]
  \centering
  \resizebox{0.86\textwidth}{!}{% Original comparison using a separable toy dataset and XOR.
\begin{tikzpicture}[>=Latex,
  axis/.style={->,thick},
  divider/.style={very thick,draw=ntuRed},
  title/.style={font=\bfseries},
  note/.style={font=\scriptsize,align=center}
]
  \begin{scope}[xshift=0cm]
    \node[title] at (2.3,3.25) {Linearly separable};
    \draw[axis] (0,0) -- (4.7,0) node[right] {$x_1$};
    \draw[axis] (0,0) -- (0,2.8) node[above] {$x_2$};
    \foreach \p in {(0.7,0.5),(1.0,1.2),(1.4,0.8),(1.7,1.5)}
      \fill[noteBlue] \p circle (3pt);
    \foreach \p in {(2.8,1.2),(3.2,1.8),(3.6,1.4),(4.0,2.2)}
      \draw[ntuRed,very thick] \p ++(-0.12,-0.12) -- ++(0.24,0.24)
        \p ++(-0.12,0.12) -- ++(0.24,-0.24);
    \draw[divider] (2.6,0.15) -- (2.15,2.65);
    \node[note] at (2.3,-0.55) {one linear boundary exists};
  \end{scope}

  \draw[dashed,gray] (5.3,-0.8) -- (5.3,3.45);

  \begin{scope}[xshift=6.2cm]
    \node[title] at (2.3,3.25) {XOR: non-linearly separable};
    \draw[axis] (0,0) -- (4.7,0) node[right] {$x_1$};
    \draw[axis] (0,0) -- (0,2.8) node[above] {$x_2$};
    \fill[noteBlue] (0.9,0.55) circle (3pt);
    \fill[noteBlue] (3.7,2.25) circle (3pt);
    \foreach \p in {(0.9,2.25),(3.7,0.55)}
      \draw[ntuRed,very thick] \p ++(-0.12,-0.12) -- ++(0.24,0.24)
        \p ++(-0.12,0.12) -- ++(0.24,-0.24);
    \node[note] at (2.3,-0.55) {no single line separates both classes};
  \end{scope}
\end{tikzpicture}
}
  \caption{Linearly separable（线性可分）数据与 XOR：后者不存在一条直线同时正确分开两类。}
  \label{fig:sc4001-l02-linear-separability}
\end{figure}

Learning rate $\alpha$ 太小会收敛缓慢，合适时稳定下降，太大则可能 oscillate（振荡）或 diverge（发散）。不存在对所有问题都最优的固定 $\alpha$；此外，讲义 GD 对 gradients 求和而 SGD 使用单样本 gradient，二者不能脱离 loss normalization（损失归一化）直接比较数值大小。SGD 的 stochasticity（随机性）有时有助于 non-convex optimization（非凸优化），但不保证得到更低 error。实践通常使用 mini-batch mean gradient（小批量平均梯度）。

\begin{figure}[htbp]
  \centering
  \resizebox{0.84\textwidth}{!}{% Original conceptual curves; not fitted to the lecture's experimental data.
\begin{tikzpicture}[>=Latex,
  axis/.style={->,thick},
  curve/.style={very thick},
  title/.style={font=\small\bfseries},
  note/.style={font=\scriptsize,align=center}
]
  \begin{scope}[xshift=0cm]
    \node[title] at (1.8,3.0) {$\alpha$ too small};
    \draw[axis] (0,0) -- (3.7,0) node[right] {iteration};
    \draw[axis] (0,0) -- (0,2.6) node[above] {$J$};
    \draw[curve,draw=noteBlue]
      plot[smooth] coordinates {(0.2,2.2)(0.8,1.85)(1.4,1.55)(2.0,1.30)(2.6,1.10)(3.3,0.93)};
    \node[note] at (1.8,-0.45) {stable but slow};
  \end{scope}

  \begin{scope}[xshift=4.5cm]
    \node[title] at (1.8,3.0) {stable $\alpha$};
    \draw[axis] (0,0) -- (3.7,0) node[right] {iteration};
    \draw[axis] (0,0) -- (0,2.6) node[above] {$J$};
    \draw[curve,draw=noteBlue]
      plot[smooth] coordinates {(0.2,2.2)(0.7,1.25)(1.2,0.68)(1.8,0.36)(2.5,0.21)(3.3,0.16)};
    \node[note] at (1.8,-0.45) {fast, monotone decrease};
  \end{scope}

  \begin{scope}[xshift=9.0cm]
    \node[title] at (1.8,3.0) {$\alpha$ too large};
    \draw[axis] (0,0) -- (3.7,0) node[right] {iteration};
    \draw[axis] (0,0) -- (0,2.6) node[above] {$J$};
    \draw[curve,draw=ntuRed]
      plot[smooth] coordinates {(0.2,2.1)(0.6,0.55)(1.0,2.25)(1.4,0.40)(1.8,2.35)(2.2,0.32)(2.6,2.45)(3.2,0.25)};
    \node[note] at (1.8,-0.45) {oscillation or divergence};
  \end{scope}
\end{tikzpicture}
}
  \caption{Learning rate 的典型影响：步长过小、稳定下降与过大振荡。曲线为概念示意，不复刻讲义实验数值。}
  \label{fig:sc4001-l02-learning-rate}
\end{figure}

PyTorch autograd（自动微分）的基本顺序是：parameters 设置 \texttt{requires\_grad=True}，调用 \texttt{loss.backward()}，在 \texttt{torch.no\_grad()} 内更新参数，再将 gradients 清为 \texttt{None}，避免默认的 gradient accumulation（梯度累积）。这只是上述 chain rule 的程序化实现。

\textbf{一分钟回顾。}
Regression 与 classification 共享 $u=\bm w^{\mathsf T}\bm x+b$，但前者预测连续值并常用 square error，后者用 sigmoid probability 与 BCE。Sigmoid--BCE 组合使 $\partial J/\partial u=p-d$；最终 label 由 $p>0.5\iff u>0$ 决定，所以 single logistic neuron 的 decision boundary 仍是 linear hyperplane。GD、SGD 与 mini-batch GD 的差异在于每次更新使用的 samples 及 gradient normalization，而不只是名称。
